\begin{exercisechunk}
\item \textbf{Repeated decisions and fresh randomness.}
\label[exercise]{ex:l18-information-budget}
\exerciseprofile{2}{2}{2}{\exproof\quad\excalculation\quad\exinterpretation}{%
Separate the number of recorded actions from the information they provide
about a teacher parameter.}
Fix $p\in[1/4,3/4]$, a horizon $H$, and a budget $m=nH$.
In family R, each teacher episode draws one
$\operatorname{Bernoulli}(p)$ bit and repeats it $H$ times.
In family I, all teacher actions are independent
$\operatorname{Bernoulli}(p)$ bits. Episodes are independent in both
families. The target in this exercise is estimation of the scalar
parameter $p$, not comparison of the two families' trajectory laws.
\begin{enumerate}[(a)]
\item Show that the un-clipped average of all $m$ recorded actions has
variance $Hp(1-p)/m$ in family R and $p(1-p)/m$ in family I.
Deduce expected absolute error bounds of $\sqrt{Hp(1-p)/m}$ and
$\sqrt{p(1-p)/m}$, respectively.
\item Suppose a separately requested initial action costs $c$ units
and gives a fresh $\operatorname{Bernoulli}(p)$ bit, where $c$ is a
positive integer and $m/c$ is an integer. A full demonstration still
costs $H$ units. In family R, compare the variances obtained by spending
all $m$ units on full demonstrations or on separate initial-action
queries. For which $c$ is the second strategy strictly better when $H>1$?
\item In family I with $c=1$, explain why choosing query stages and
prefixes adaptively cannot produce more informative observations than
the $m$ independent actions already obtained passively. Assume the
teacher uses fresh independent randomness for every response.
\end{enumerate}
\end{exercisechunk}
